\section*{Chapter \ref{ch:hf:short-horizon-alpha} --- Short-Horizon Alpha}

\begin{solution}{exo:hf:short-horizon-alpha:1}
Imbalance $(12-4)/16=0.5$; forecast $0.30\times0.5=0.15$ ticks.
\end{solution}

\begin{solution}{exo:hf:short-horizon-alpha:2}
Half a tick plus the fee, 0.1 cent $=0.1$ tick: 0.6 ticks.
\end{solution}

\begin{solution}{exo:hf:short-horizon-alpha:3}
Bid up: $+3$ (the new bid size, nothing removed); ask unchanged in price: $-6+5=-1$. The increment is $3-1=2$ lots.
\end{solution}

\begin{solution}{exo:hf:short-horizon-alpha:4}
The mid lags the efficient price for seconds, so a recent move tends to continue while the book catches up. In a real stock, where quotes follow
information faster and impact partly reverts, the own move usually enters with a small or negative coefficient at these horizons.
\end{solution}

\begin{solution}{exo:hf:short-horizon-alpha:5}
Imbalance (0.187) at half a second; order-flow imbalance (0.372) at ten. A quoter holding for five seconds should weight the slower features:
fit the forecast at the horizon it holds, not the shortest one.
\end{solution}

\begin{solution}{exo:hf:short-horizon-alpha:6}
The forecast's typical size is its correlation times the target's volatility, much smaller than the target itself; only a tiny share of forecasts
exceed the cost of crossing (0.01\% exceeded 0.9 ticks here). Its value lies in many small decisions about which quote to show.
\end{solution}

\begin{solution}{exo:hf:short-horizon-alpha:7}
0.444, with ten messages taking a median of 451 milliseconds.
\end{solution}

\begin{solution}{exo:hf:short-horizon-alpha:8}
It executes at the mid, at the instant of the signal: taking costs half the spread and the fee, and the price available a message later has already
moved with the signal. Rerun at the ask (or bid) of the next message with fees; the edge per trade falls by at least the half-spread.
\end{solution}

\begin{solution}{pb:hf:short-horizon-alpha:1}
\begin{enumerate}
\item A forecast over milliseconds to seconds from book, trades and related instruments, used mainly to decide which quotes to show; intraday alpha
works at minutes and pays for crossing.
\item Queue imbalance; order-flow imbalance over one second; net aggressive volume over one second; the leader's mid change over one second; the own
mid change over one second.
\item 0.320, 0.285, 0.077, 0.288, 0.150; combined 0.480.
\item A fully bid-imbalanced book forecasts $+0.30$ ticks; a tick of the leader's move $+0.34$ ticks.
\item The simulated mid trails its efficient price for seconds (lag-one correlation of two-second returns 0.44), which makes the next move unusually
predictable.
\item 0.28 at half a second, 0.37 at one, 0.48 at two, 0.54 at five, 0.50 at ten.
\item About two average price changes.
\item 0.477 (33 ms), 0.405 (0.9 s), 0.303 (2.3 s).
\item Mean realised changes rise with the forecast across tenths, matching it closely at the extremes (top tenth 0.31 forecast, 0.33 realised).
\item Moving or withdrawing the side the forecast says will be picked off; in one tick: showing only the favourable side.
\item Cross when the forecast exceeds half the spread plus the taker fee plus a margin: $0.5+0.1+0.3=0.9$ ticks.
\item 0.01\% of forecasts beyond 0.9 ticks, 5.8\% beyond 0.3.
\item None 16\,550 at 0.294; skew 11\,050 at 0.451; skew and take 10\,867 at 0.452; one message late 10\,883 at 0.414.
\item 0.451 against 0.294 ticks a share; one message late the passive mark-out falls by 8\% (0.452 to 0.414), while the forecast's information
coefficient falls by less than 1\% (0.480 to 0.477).
\item About 25 takes in six sessions: their mark-out has a standard error as large as the differences.
\item Both withdraw a side; fading uses a forecast of the fill's mark-out, the skew a forecast of the price's direction; both give up a third of the
volume for better fills.
\item Post-sweep reversion: in this market the own move predicts continuation.
\item Median arbitrage life fell from 97 ms to 7 ms while the median profit per opportunity stayed near 0.08 index points.
\item Execute a message later at the then-available price with fees, and count enough takes (more sessions) for a standard error.
\item Deciding which of its own quotes to leave in the book.
\end{enumerate}
\end{solution}

\begin{solution}{iq:hf:short-horizon-alpha:1}
Withdraw or raise the fund's ask: the future's move will reach the fund, and a stale ask will be lifted by whoever reads the future faster; keep the
bid.
\iqlookfor{the leader's move as a forecast and protecting the stale side.}
\end{solution}

\begin{solution}{iq:hf:short-horizon-alpha:2}
The change of the bid queue (counting a better bid as its new size, a worse one as the loss of the old) minus the same for the ask, summed over
events. It includes cancellations and new limit orders, which move the price as much as trades do, while signed volume sees only executions.
\iqlookfor{the definition and why cancellations matter.}
\end{solution}

\begin{solution}{iq:hf:short-horizon-alpha:3}
Overlapping targets make neighbouring errors correlated: use HAC standard errors (One Quant Book 7, chapter 6), evaluate on sessions not used for
fitting, and report the information coefficient by horizon.
\iqlookfor{overlap, HAC errors, out-of-sample sessions.}
\end{solution}

\begin{solution}{iq:hf:short-horizon-alpha:4}
One feature engine for both, fed message by message; a replay test that runs the production engine over stored feeds and compares its output to
research's, value by value.
\iqlookfor{a single implementation and a parity test.}
\end{solution}

\begin{solution}{iq:hf:short-horizon-alpha:5}
When the forecast exceeds half the spread plus the taker fee plus a margin; the margin covers the forecast's error, its age by the time the order
arrives, and the chance that faster traders took the move first.
\iqlookfor{the threshold and its components.}
\end{solution}

\begin{solution}{iq:hf:short-horizon-alpha:6}
The conditional mean has standard deviation $\rho s$, so $P(|\hat\alpha|>c)=2\Phi(-c/(\rho s))$: with $\rho s=0.288$ and $c=0.9$, $2\Phi(-3.13)
\approx0.18\%$.
\iqlookfor{the variance of a forecast and a Gaussian tail.}
\end{solution}
